\section{Introduction}
The $n$-queens problem asks how many ways there are to place $n$ queens on an $n\times n$ board so
that no two attack each other. We denote this number by $Q(n)$ and count solutions related by a
symmetry separately. We number rows and columns starting at one, with rows running from top to
bottom and columns from left to right, and $(i,j)$ is the square in row $i$ and column $j$. Reading
the columns of the queens in row order identifies each solution with a permutation $\sigma$, whose
queens occupy the squares $(i,\sigma(i))$. The diagonal constraints require the differences
$i-\sigma(i)$ to be distinct, and likewise the sums $i+\sigma(i)$; they thus determine which of the
$n!$ permutations qualify. For example, $(2,4,1,3)$ and $(3,1,4,2)$ are the two solutions on the
$4\times4$ board.

To study how $Q(n)$ grows, it is convenient to take logarithms. Earlier combinatorial results
include the upper bound of Luria~\cite{luria}, the lower bound of Luria and Simkin~\cite{ls}, and
the work of Bowtell and Keevash~\cite{bk}. Bell and Stevens~\cite{bs} give an overview.
Simkin's variational formula~\cite{simkin} determines $\log Q(n)$ up to an $o(n)$ error.
In the normalization of the variational profile~\cite{pv}, the leading term is $\log(n!)+n\Phi$.
Nobel, Agrawal, and Boyd~\cite{nab} obtained numerical bounds for the first-order constant.
The profile describes the distribution of mass on the board and along each diagonal. These data
enter the canonical predictor $Q_{\rm pred}(n)$ of~\cite{pc}. An entropy is calculated from the
square weights and a correction is added for the dependence between the lines.
Appendix~\ref{app:predictor} gives the finite formula that we use throughout the paper.

The next advance was to identify the constant term. After subtracting $\log(n!)+n\Phi$, the
logarithm of the predictor tends to $C_{\rm pred}$~\cite{pc}. The exact count has that same
constant~\cite{tc}. Thus, on forming the signed difference
\[
 R_Q(n)=\log Q(n)-\log Q_{\rm pred}(n),
\]
both the leading term and the constant cancel. A preliminary step in the present work, using
magnitude estimates, gives $|R_Q(n)|\le C/n$ for sufficiently large $n$. The proof traceability
table in the supplement identifies the source of this bound. Appendix~\ref{app:signed} uses it in
developing the stopped-history identity. Thus $nR_Q(n)$ is bounded, although it could still
oscillate with the size of the board. Does it converge? Theorem~\ref{teo:principal} identifies
this limit, which we denote by \(\cerr\). While \(C_{\rm pred}\) is the constant term shared by the
two logarithms, \(\cerr/n\) is the first relative correction in passing from the finite predictor
to the exact count.

The difficulty can be seen by revealing a solution row by row. While many rows remain, each
insertion introduces a small defect, and it is their signed sum that matters. When few rows remain,
two squares that are still available may share a diagonal. The completability results of~\cite{ef}
describe the role of the final rows; Glock, Munhá Correia, and Sudakov~\cite{gms} studied when a
partial placement can be completed. To calculate the coefficient we also need the weight of these
conflicts in the error. Following the early cut-off of~\cite{ef}, we stop the process when a fixed
fraction of rows remains empty. This allows us to study the defect accumulated before the cut-off
and to measure the error on the residual board. Comparing two cut-offs will bring the two
contributions together. Figure~\ref{fig:mascara-ilustrada} follows the first four-queens solution:
revealing a single queen already shows the difference between a rook completion and a queen
completion.
\begin{figure}[htbp]
\centering
\begin{tikzpicture}[x=.61cm,y=.61cm]
  \bcuadricula{4}
  \foreach \r/\c in {1/2,2/4,3/1,4/3}{\breina{4}{\r}{\c}}
  \node[font=\small] at (2,-.62) {Complete solution};
  \begin{scope}[shift={(7,0)}]
    \bcuadricula{4}
    \fill[orange,opacity=.24] (0,3) rectangle (4,4);
    \fill[blue,opacity=.15] (1,0) rectangle (2,4);
    \begin{scope}
      \clip (0,0) rectangle (4,4);
      \draw[teal!75!black,line width=.7pt] (-2,7)--(6,-1);
      \draw[red!70!black,dashed,line width=.7pt] (-2,0)--(6,8);
    \end{scope}
    \foreach \r/\c in {2/4,3/1,3/3,4/1,4/3,4/4}{\bmarca{4}{\r}{\c}}
    \breina{4}{1}{2}
    \node[font=\small] at (2,-.62) {One queen revealed};
  \end{scope}
  \begin{scope}[shift={(14,.1)}]
    \foreach \r/\cells in {1/{0,0,1},2/{1,1,0},3/{1,1,1}}{
      \foreach \w [count=\c] in \cells {
        \ifnum\w=1 \fill[teal!16] (\c-1,3-\r) rectangle (\c,4-\r);\fi
        \node[font=\small] at (\c-.5,3.5-\r) {$\w$};
      }
    }
    \draw[step=1,black!45,line width=.3pt] (0,0) grid (3,3);
    \draw[btinta,line width=.6pt] (0,0) rectangle (3,3);
    \foreach \i/\r in {1/2,2/3,3/4}{\node[font=\scriptsize] at (-.3,3.5-\i) {$\r$};}
    \foreach \i/\c in {1/1,2/3,3/4}{\node[font=\scriptsize] at (\i-.5,3.32) {$\c$};}
    \node[font=\small] at (1.5,-.72) {Residual mask};
  \end{scope}
\end{tikzpicture}
\caption{From a solution to the rook problem. After $(1,2)$ is revealed, the dots in the middle are the
ones in the mask on the right, shown with their original labels. Row 2 requires column 4. Columns 1 and 3
can be assigned to rows 3 and 4 in two ways: $\operatorname{per}M=2$. Only the assignment on the left
completes the queen placement. In the other, $(2,4)$ and $(3,3)$ attack each other. Thus $Q(B)=1$.}
\label{fig:mascara-ilustrada}
\end{figure}

